% Staged section: include from extremal.tex only after the complete enumeration
% replay and physical-classification axiom audit have passed.
\subsection{The complete fourteen-cube classification}\label{sec:classification14}
The preceding family accounts for most, but not all, extremal obstructions.
There are two further forms, both using six varieties. Throughout this
subsection equivalence means common colour relabelling and optional
simultaneous reflection of the collection. A labelled collection keeps the
five colour names fixed; identical physical copies remain indistinguishable.

\paragraph{An all-adjacent family.}
Choose a colour $o$, an ordered partition $P,Q$ of $V\setminus\{o\}$ into
pairs, and a derangement $\rho$ of those four colours. Take one copy of each
repeat-$o$ variety with cuts $C(P)$ and $C(P\cup\{o\})$. For each $v\in P$,
take three copies with cut $C(\{v,o\})$ and repeat $\rho(v)$; for each
$v\in Q$, take three copies with cut $C(V\setminus\{v,o\})$ and repeat
$\rho(v)$. The condition $\rho(v)\ne v$ makes every variety well defined.
These collections have multiplicities $(3,3,3,3,1,1)$.

Their tricoloured union omits exactly $Q\times P$. Its bidirected graph is
the two triangles on $P\cup\{o\}$ and $Q\cup\{o\}$, sharing only $o$.
It has no four-cycle, so no opposite-repeat target is covered. For any
repeat other than $o$, the bicoloured corners fix the smaller pair of its
unique three-copy variety; the reverse cut requires a missing arc.
For repeat $o$, a smaller pair crossing $P,Q$ also requires a missing arc
in either orientation. Thus precisely the six source varieties are covered.

For $v\in P$, with $w$ the other member of $P$, the three-copy target's
private arcs are $vw,ow$. For $v\in Q$, with $w$ the other member of $Q$,
they are $wv,wo$. Its two bicoloured corners are also private, since $\rho$
is bijective. Each three-copy variety therefore faces four private positions, and each
singleton faces its two private bicoloured positions. Every target fails.

The colour $o$ is recovered from the two singleton varieties, and the
missing cut recovers $P,Q$ up to reversal. The eight-element partition
stabilizer acts on $\rho$ by conjugation. For $o=4$, $P=\{0,3\}$ and
$Q=\{1,2\}$, its four orbits are represented by
\[
\begin{array}{c|c|r}
 \text{type}&\rho&\text{labelled orbit size}\\\hline
 \text{within-pair transpositions}&(0\ 3)(1\ 2)&30\\
 \text{cross-pair transpositions}&(0\ 1)(3\ 2)&60\\
 \text{alternating four-cycle}&(0\ 1\ 3\ 2)&60\\
 \text{nonalternating four-cycle}&(0\ 3\ 1\ 2)&120
\end{array}
\]
The normalized orbit sizes are $1,2,2,4$, accounting for the nine
derangements. There are $5\binom42=30$ ordered templates, giving 270
labelled collections in these four classes.

\paragraph{A mixed exception.}
The remaining form is the following single collection, in the face order
$(+x,-x,+y,-y,+z,-z)$:
\[
\begin{array}{c|c|c|c}
 \text{variety}&\text{face tuple}&\text{copies}&\text{private positions}\\\hline
 X_0&(0,1,0,2,3,4)&5&6\\
 X_1&(0,4,1,2,3,1)&1&2\\
 X_2&(0,4,1,2,3,2)&1&2\\
 X_3&(0,4,1,3,3,2)&1&2\\
 O&(0,4,1,2,3,3)&3&4\\
 Y&(0,4,1,4,3,2)&3&4
\end{array}
\]
Only $O$ is opposite-repeat. The arc union misses $04,21,23,34$; its
bidirected edges are $01,02,03,13,14,24$. The only four-cycle is
$0\!-\!1\!-\!4\!-\!2\!-\!0$, corresponding to $O$. For each repeated colour,
exactly two bicoloured corners occur. They fix the listed adjacent variety
up to reversal, and every reverse requires a missing arc. These six
varieties are therefore exactly the covered targets. The private
tricoloured positions are $30,31,32,40$ for $X_0$, none for $X_1,X_2,X_3$,
$10,14,20,24$ for $O$, and $12,13$ for $Y$; every adjacent target also has
two private bicoloured positions. Each target needs one more copy than is
available. The collection fails and its labelled orbit has size 240.

\begin{theorem}[Fourteen-cube obstructions]\label{thm:classification14}
Every fourteen-cube obstruction is equivalent to a member of exactly one
of the following forms:
\begin{enumerate}
 \item the failing collections of Corollary~\ref{cor:extremal-family};
 \item the all-adjacent derangement family above;
 \item the mixed exception above.
\end{enumerate}
There are respectively $51,4,1$ symmetry classes, hence 56 in total.
Their distribution by number of varieties is
\[
\begin{array}{c|rrrr|r}
 \text{varieties}&3&4&5&6&\text{total}\\\hline
 \text{classes}&5&20&26&5&56\\
 \text{labelled collections}&840&4800&5280&510&11430.
\end{array}
\]
\end{theorem}
\begin{proof}
The preceding arguments prove failure of all three forms. In the first
form, there are 91 distributions satisfying the disjoint-repeat condition.
The only residual nonidentity symmetry sends
\[
 (x_2,x_3,x_4;y_1,y_2,y_4)
 \longmapsto (y_4,y_1,y_2;x_3,x_4,x_2).
\]
It fixes 11 distributions: $y_4=x_2$, $y_1=x_3$, $y_2=x_4$ and
$x_2x_4=0$. Burnside's lemma gives $(91+11)/2=51$ classes. The three
forms cannot be equivalent, since their opposite-repeat totals are
respectively four, zero and three.

Necessity is the finite part of the proof. The anchored Hall cover applies
to fourteen sources. Theorem~\ref{thm:multiplicity-five}, independently of
the fifteen-cube threshold, removes the eight roots with anchor
multiplicity six or seven. The remaining 349 roots use total fourteen.
Their interval certificates retain 81 explicit terminal count vectors.
In addition to the exclusion, refinement and splitting rules, the checker
accepts an exact singleton box only when its vector occurs in the terminal
list. Thus every admissible vector in every root belongs to that list.
Each terminal is carried to one of the 56 representatives by an explicit
common colour/reflection symmetry. Separate finite checks verify the
Hall deficiencies, the full symmetry group and pairwise inequivalence.
The retained data give the displayed orbit and support counts.
\end{proof}

In particular, any fourteen-cube collection using at least seven varieties solves. The all-adjacent obstructions
are exactly the derangement family.

The classification also explains the corner threshold by a short deletion
argument. Suppose an unsolvable fifteen-cube pool existed. Every deletion
would be a classified fourteen-cube obstruction. If all sources were
adjacent-repeat, every deletion would have profile $(3,3,3,3,1,1)$.
A variety occurring three times in one deletion had multiplicity three or four in the
original pool. Deleting from that variety in the first case, or from
another variety in the second, leaves multiplicity two or four, both
impossible.

Otherwise there is exactly one opposite-repeat variety $A$: two distinct
ones could both be retained by a deletion. There are adjacent sources as
well, since every classified obstruction has at most four opposite copies.
Deleting an adjacent copy shows that the multiplicity of $A$ is three or
four; deleting an $A$ shows that one less is also three or four. Hence it
is four. Deleting an $A$ gives the mixed exception, so the adjacent
multiplicities in the original pool are $(5,3,1,1,1)$. Delete from the
five-copy variety. The resulting obstruction has four opposite copies
and six varieties, whereas the first form has at most five varieties.
This final contradiction proves the fifteen-cube upper bound.
